\documentclass[12pt, a4paper]{article}

% === PACKAGES ===
\usepackage[utf8]{inputenc}
\usepackage[T1]{fontenc}
\usepackage[english]{babel}
\usepackage{amsmath, amssymb, amsthm}
\usepackage{geometry}
\usepackage{hyperref}
\usepackage{xcolor}
\usepackage{booktabs}
\usepackage{float}
\usepackage{tcolorbox}
\usepackage{fancyhdr}
\usepackage{lastpage}
\usepackage{graphicx}

% === GEOMETRY ===
\geometry{top=2.5cm, bottom=2.5cm, left=2.5cm, right=2.5cm}

% === HYPERREF ===
\hypersetup{
    colorlinks=true,
    linkcolor=blue!70!black,
    urlcolor=cyan!70!black,
    citecolor=red!70!black,
    pdftitle={Ontological Density Field v3.1},
    pdfauthor={Fedor Kapitanov}
}

% === THEOREMS ===
\newtheorem{definition}{Definition}
\newtheorem{proposition}{Proposition}
\newtheorem{theorem}{Theorem}
\newtheorem{lemma}{Lemma}
\newtheorem{corollary}{Corollary}
\newtheorem{remark}{Remark}
\newtheorem{axiom}{Axiom}
\newtheorem{conjecture}{Conjecture}

% === HEADERS ===
\pagestyle{fancy}
\fancyhf{}
\fancyhead[L]{\small Ontological Density Field $I(x)$ — v3.1}
\fancyhead[R]{\small F. Kapitanov}
\fancyfoot[C]{\thepage\ / \pageref{LastPage}}
\renewcommand{\headrulewidth}{0.4pt}
\renewcommand{\footrulewidth}{0.4pt}

% === TCOLORBOX ===
\tcbuselibrary{skins, breakable}

% === DOCUMENT ===
\begin{document}

% ==============================================================================
% TITLE PAGE
% ==============================================================================

\begin{titlepage}
\centering
\vspace*{1cm}

{\Huge\bfseries Ontological Density Field $I(x)$}

\vspace{0.6cm}

{\Large\itshape A Deductive Model of Spacetime Information Content}

\vspace{0.4cm}

{\large Implications for $H_0$ Tension, Dark Matter, and the $(\pi-3)$ Spectral Signature}

\vspace{0.3cm}

{\large\bfseries Version 3.1 — Spectral Gap Extension}

\vspace{1.5cm}

{\Large\bfseries Fedor Kapitanov}

\vspace{0.5cm}

{\large Independent Researcher, Moscow}\\[0.3cm]
{\large ORCID: \href{https://orcid.org/0009-0009-6438-8730}{0009-0009-6438-8730}}\\[0.3cm]
{\large \href{mailto:prtyboom@gmail.com}{\texttt{prtyboom@gmail.com}}}

\vspace{1.5cm}

{\large \today}

\vspace{1cm}

\begin{tcolorbox}[colback=blue!5!white, colframe=blue!75!black, width=0.95\textwidth]
\textbf{Abstract}

\vspace{0.2cm}

We present a deductive framework in which the finiteness and discreteness of observable reality follow from set-theoretic considerations regarding the observer's capacity. By formalizing observation as a mapping into a finite internal state space, we prove that the set of operationally distinguishable states is finite, regardless of the cardinality of the underlying "Absolute".

This theorem provides the logical foundation for an \emph{ontological density field} $I(x)$, representing the density of distinguishable phase-space cells. Physical constants ($\hbar$, $c$, $G$) are interpreted as encoding parameters relating discrete information to continuous geometry.

\textbf{New in v3.1}: We demonstrate that the spectral action on a discrete $N$-cycle graph $C_N$ differs from the continuous circle $S^1$ by a defect proportional to $(\pi - 3)$. For $N=12$ (zodiacal discretization), numerical computation yields:
\begin{equation*}
    |\Delta S| = 52 \times (\pi - 3) \quad \text{with 99.47\% accuracy.}
\end{equation*}
This provides numerical evidence that $(\pi - 3) \approx 0.14159$ emerges as a fundamental signature of discrete-to-continuum transition.

\end{tcolorbox}

\vfill

{\small This work is licensed under CC BY 4.0}

\end{titlepage}

\tableofcontents
\newpage

% ==============================================================================
% SECTION 1: INTRODUCTION
% ==============================================================================

\section{Introduction}

\subsection{Motivation}

The standard approach to fundamental physics takes quantum mechanics and general relativity as foundational, then derives information-theoretic results (Bekenstein bound, holographic principle) as consequences. This creates a logical structure where the finiteness of observable reality depends on specific physical theories.

We invert this hierarchy. We show that finiteness follows from the logical structure of observation itself. If an observer has finite capacity, the observable world must be finite and discrete. Physical theories then emerge as descriptions of \emph{how} this finite structure is organized.

\subsection{Structure of the Paper}

\begin{itemize}
    \item \textbf{Section 2}: Axiomatic foundation and the Finiteness Theorem via equivalence classes.
    \item \textbf{Section 3}: Interpretation of physical constants as conversion factors.
    \item \textbf{Section 4}: Operational definitions and the ontological density field.
    \item \textbf{Section 5}: Covariant action and field equations.
    \item \textbf{Sections 6--9}: Applications to Dark Matter, Hubble Tension, and predictions.
    \item \textbf{Section 10}: Spectral gap and the $(\pi-3)$ signature (new in v3.1).
\end{itemize}

% ==============================================================================
% SECTION 2: AXIOMATIC FOUNDATION
% ==============================================================================

\section{Axiomatic Foundation}

\subsection{The Three Axioms}

We begin with three axioms that make no reference to physics.

\begin{axiom}[Totality]
There exists a set $\Omega$ of all possible states (the ``Absolute''). The cardinality of $\Omega$ is unspecified and may be infinite.
\end{axiom}

\begin{axiom}[Distinction]
There exists an observer $O$ that is a proper subset of $\Omega$:
\begin{equation}
    O \subset \Omega, \quad O \neq \Omega.
\end{equation}
The observer is defined by its separation from the remainder $R = \Omega \setminus O$.
\end{axiom}

\begin{axiom}[Operational Existence]
A state $r \in R$ exists for observer $O$ if and only if $O$ can distinguish $r$ from other states. Indistinguishable states are operationally identical.
\end{axiom}

\begin{remark}
These axioms are minimal. Axiom~1 posits that ``something exists.'' Axiom~2 posits that ``distinction occurs.'' Axiom~3 defines existence through distinguishability --- the standard operationalist criterion used throughout physics.
\end{remark}

\subsection{Observable Equivalence Classes}

Observation is formalized as a map
\begin{equation}
    f: R \to O,
\end{equation}
assigning to each external state $r \in R$ an internal state $f(r) \in O$ of the observer.

\begin{definition}[Observational Equivalence]
Given an observation map $f: R \to O$, we define an equivalence relation $\sim$ on $R$ by
\begin{equation}
    r_1 \sim r_2 \;\;\Longleftrightarrow\;\; f(r_1) = f(r_2).
\end{equation}
The set of \emph{operationally distinguishable} states is the quotient set
\begin{equation}
    R/\!\sim \;=\; \{[r] \mid r \in R\}, \quad [r] = \{ r' \in R \mid r' \sim r \}.
\end{equation}
\end{definition}

By construction, there is a natural bijection
\begin{equation}
    R/\!\sim \;\simeq\; \mathrm{im}\,f \subseteq O,
\end{equation}
so that the number of distinguishable states is
\begin{equation}
    N_{\text{obs}} = |R/\!\sim| = |\mathrm{im}\,f|.
\end{equation}

\subsection{Finiteness of Observable Reality}

\begin{theorem}[Finiteness of Observable Reality]\label{thm:finiteness}
Let $\Omega$ be a set of states, $O \subset \Omega$ an observer, and $R = \Omega \setminus O$ the remainder. Assume that the internal state space of the observer is finite:
\begin{equation}
    |O| = N_O < \infty.
\end{equation}
Then the number of states in $R$ that are operationally distinguishable by $O$ is finite, and satisfies
\begin{equation}
    N_{\text{obs}} \;=\; |R/\!\sim| \;=\; |\mathrm{im}\,f| \;\le\; |O| = N_O < \infty.
\end{equation}
\end{theorem}

\begin{proof}
\textbf{Step 1.} By assumption, the internal state space of the observer is finite: $|O| = N_O < \infty$.

\textbf{Step 2.} Any observable state of $R$ corresponds to a distinct internal state of $O$, i.e.\ to an element of $\mathrm{im}\,f \subseteq O$.

\textbf{Step 3.} Therefore the number of operationally distinguishable states is
\begin{equation}
    N_{\text{obs}} = |\mathrm{im}\,f| \le |O| = N_O < \infty.
\end{equation}
\end{proof}

\begin{corollary}[Operational Discreteness]\label{cor:discrete}
For any observer with finite internal state space, the set $R/\!\sim$ of operationally distinguishable states is finite and therefore admits only a discrete topology. Any continuous description of reality is, for such an observer, an idealization.
\end{corollary}

% ==============================================================================
% SECTION 3: EMERGENCE OF PHYSICAL CONSTANTS
% ==============================================================================

\section{Emergence of Physical Constants}

Given that observable reality, for any finite observer, consists of a finite number $N$ of distinguishable states (Theorem~\ref{thm:finiteness}), we now reinterpret the familiar physical constants as encoding parameters relating this discrete structure to continuum quantities.

\subsection{The Counting Unit: $\hbar$}

\begin{definition}[Planck Constant as Conversion Factor]
In this framework, the reduced Planck constant $\hbar$ plays the role of a conversion factor between the number of distinguishable states and phase-space volume. We write
\begin{equation}
    V_{\text{cell}} = \hbar^3
\end{equation}
for the phase-space volume associated with one unit of distinguishability.
\end{definition}

\subsection{The Update Rate: $c$}

\begin{definition}[Speed of Light as Causality Limit]
The speed of light $c$ is interpreted as the maximal rate at which distinctions can propagate through the underlying causal graph.
\end{definition}

\subsection{The Packing Coefficient: $G$}

\begin{definition}[Maximal Entropy--Area Relation]
We postulate that the maximal entropy $S_{\max}$ associated with a bounded region of spacetime is proportional to the area $A$ of an appropriate bounding surface:
\begin{equation}
    S_{\max} = \eta\, A, \quad \eta = \frac{c^3}{4 G \hbar}.
\end{equation}
\end{definition}

% ==============================================================================
% SECTION 4: OPERATIONAL DEFINITIONS
% ==============================================================================

\section{Operational Definitions}

\subsection{Ontological Density}

\begin{definition}[Ontological Density Field]
The ontological density at spacetime point $x$ is defined as the density of distinguishable states per unit proper volume:
\begin{equation}
    I(x) \equiv \frac{N(x)}{V_{\text{proper}}} = \frac{1}{\hbar^3} \int \frac{d^3p}{(2\pi)^3} f(x,p).
\end{equation}
\end{definition}

\subsection{Distinction Creation Rate}

\begin{definition}[Source Term]
The rate of distinction creation is:
\begin{equation}
    S_I(x) = \frac{\rho(x) c^2}{\hbar \tau_*} = \frac{\rho(x) c^2 k_B T}{\hbar^2},
\end{equation}
where $\tau_* = \hbar/(k_B T)$ is the thermal relaxation time.
\end{definition}

% ==============================================================================
% SECTION 5: COVARIANT ACTION AND FIELD EQUATIONS
% ==============================================================================

\section{Covariant Action and Field Equations}

\subsection{Action}

The total action is:
\begin{equation}
    S = \frac{c^4}{16\pi G} \int d^4x \sqrt{-g} \, R + \int d^4x \sqrt{-g} \, \mathcal{L}_I + S_{\text{matter}},
\end{equation}
with:
\begin{equation}
    \mathcal{L}_I = -\frac{1}{2} f(I) g^{\mu\nu} \nabla_\mu I \nabla_\nu I - V(I) + \kappa_I I T.
\end{equation}

\subsection{Kinetic Function}

To ensure causality and correct dimensions:
\begin{equation}
    f(I) = \frac{\hbar c^3}{3 I_0^{1/2}} \left(\frac{I}{I_0}\right)^{-1/2}.
\end{equation}

\subsection{Field Equation}

Variation with respect to $I$ yields:
\begin{equation}
    \nabla_\mu \left( f(I) \nabla^\mu I \right) - \frac{1}{2} f'(I) (\nabla I)^2 + V'(I) = \kappa_I T.
\end{equation}

% ==============================================================================
% SECTION 6: SCREENING MECHANISM
% ==============================================================================

\section{Screening Mechanism}

\subsection{Effective Potential}

\begin{equation}
    V_{\text{eff}}(I, \rho) = V_0 \left(\frac{I_{\text{screen}}}{I}\right)^2 + \kappa_I I \rho c^2.
\end{equation}

\subsection{Effective Coupling}

\begin{equation}
    \alpha_{\text{eff}}(\rho) = \frac{\alpha_0}{1 + (\rho/\rho_{\text{screen}})^{2/3}}.
\end{equation}

% ==============================================================================
% SECTION 7: GHOST INFORMATION AND DARK MATTER
% ==============================================================================

\section{Ghost Information and Dark Matter Analogs}

\subsection{Accumulation}

\begin{equation}
    I_{\text{ghost}}(x,t) = \int_0^t S_I(x,t') e^{-\Lambda_I(t-t')} dt'.
\end{equation}

\subsection{Effective Dark Matter Density}

\begin{equation}
    \rho_{\text{ghost}}(r) = \frac{\kappa}{c^2} \left(I_{\text{ghost}}(r) - I_0\right).
\end{equation}

% ==============================================================================
% SECTION 8: HUBBLE PARAMETER VARIATIONS
% ==============================================================================

\section{Hubble Parameter Variations}

\subsection{Effective Hubble Parameter}

\begin{equation}
    H_0^{\text{obs}} \approx H_0^{\text{true}} \left[1 + \frac{\alpha_0}{2} \ln\left(\frac{I}{I_0}\right)\right].
\end{equation}

\subsection{Environment Dependence}

\begin{center}
\begin{tabular}{lcc}
\toprule
\textbf{Environment} & $I/I_0$ & $H_0^{\text{obs}}$ (km/s/Mpc) \\
\midrule
Cosmic void & $0.7$ & $65 \pm 1$ \\
CMB (cosmic mean) & $1.0$ & $67.4 \pm 0.5$ \\
Local Sheet & $1.5$ & $69 \pm 1$ \\
Galaxy host & $3.0$ & $73 \pm 1.5$ \\
\bottomrule
\end{tabular}
\end{center}

% ==============================================================================
% SECTION 9: FALSIFIABLE PREDICTIONS
% ==============================================================================

\section{Falsifiable Predictions}

\begin{enumerate}
    \item \textbf{Environment--$H_0$ correlation}: Systematic variation of $H_0$ with local density.
    \item \textbf{Core scaling relation}: $R_{\text{core}} \propto M_{\text{halo}}^{1/3}$.
    \item \textbf{Scalar GW polarization}: Breathing-mode with $v_I/c \approx 0.58$.
\end{enumerate}

% ==============================================================================
% SECTION 10: SPECTRAL GAP AND THE (π-3) SIGNATURE
% ==============================================================================

\section{Spectral Gap and the $(\pi - 3)$ Signature}

\subsection{Motivation}

The transition from continuous to discrete geometry must leave a measurable imprint. We investigate this via the spectral action --- the trace of the heat kernel of the Dirac operator.

\subsection{Setup}

We compare:
\begin{itemize}
    \item \textbf{Continuous}: Circle $S^1$ of circumference $L = 2\pi$
    \item \textbf{Discrete}: Cyclic graph $C_N$ with $N$ vertices
\end{itemize}

The discrete Dirac operator with antiperiodic (spinorial) boundary conditions:
\begin{equation}
    D_N = \frac{1}{2a}\begin{pmatrix} 0 & \nabla \\ \nabla^\dagger & 0 \end{pmatrix},
\end{equation}
where $a = 2\pi/N$ is the lattice spacing and $\nabla$ implements the twisted shift with $\nabla_{N-1,0} = -1$ (antiperiodicity).

The spectral action with UV cutoff $\Lambda$:
\begin{equation}
    S[\Lambda] = \mathrm{Tr}\left( e^{-D^2/\Lambda^2} \right).
\end{equation}

\subsection{Numerical Result}

The spectral defect is defined as:
\begin{equation}
    \Delta S = S_{\text{cont}} - S_{\text{disc}}.
\end{equation}

For $N = 12$ and $\Lambda = 10$, numerical diagonalization yields:

\begin{tcolorbox}[colback=yellow!10!white, colframe=orange!75!black, title=Key Numerical Result]
\begin{equation}
    |\Delta S| = 7.3236
\end{equation}
\begin{equation}
    \frac{|\Delta S|}{52} = 0.14084
\end{equation}
\begin{equation}
    \pi - 3 = 0.14159
\end{equation}
\textbf{Match accuracy: 99.47\%}
\end{tcolorbox}

\subsection{The 52-fold Structure}

We find:
\begin{equation}
    \boxed{|\Delta S(N=12, \Lambda=10)| = 52 \times (\pi - 3)}
\end{equation}

The factor 52 admits the factorization:
\begin{equation}
    52 = 4 \times 13,
\end{equation}
where 4 corresponds to spinorial degrees of freedom and 13 is prime.

\subsection{Divisor Analysis}

Among nearby integers, 52 provides the optimal match:

\begin{center}
\begin{tabular}{cc}
\toprule
Divisor & Match \% \\
\midrule
50 & 96.6 \\
51 & 98.6 \\
\textbf{52} & \textbf{99.5} \\
53 & 97.6 \\
54 & 95.8 \\
\bottomrule
\end{tabular}
\end{center}

\subsection{Parameter Dependence}

The match is optimal at $\Lambda = 10$:

\begin{center}
\begin{tabular}{ccc}
\toprule
$\Lambda$ & $|\Delta S|/52$ & Match \% \\
\midrule
1 & 0.107 & 75.7 \\
5 & 0.252 & 21.9 \\
\textbf{10} & \textbf{0.141} & \textbf{99.5} \\
20 & 0.047 & 33.3 \\
\bottomrule
\end{tabular}
\end{center}

\subsection{Interpretation}

The number $N = 12$ represents the minimal ``zodiacal'' discretization of the circle --- the simplest division compatible with both 3-fold and 4-fold symmetry.

The appearance of $(\pi - 3)$ in the spectral defect suggests that this irrational quantity encodes the fundamental mismatch between continuous and discrete geometry:
\begin{equation}
    \pi - 3 = 0.14159\ldots
\end{equation}

This is the ``leftover'' of $\pi$ after subtracting the largest integer less than $\pi$. In ORT terms, it represents the irreducible information loss when discretizing a circle.

\subsection{Conjecture}

\begin{conjecture}[$(\pi-3)$ Universality]
The quantity $(\pi - 3)$ appears universally in the spectral defect between continuous manifolds and their minimal discrete approximations, with integer prefactors encoding topological and spinorial data.
\end{conjecture}

% ==============================================================================
% REFERENCES
% ==============================================================================

\newpage
\begin{thebibliography}{99}

\bibitem{Cantor1891}
G.~Cantor,
``\"Uber eine elementare Frage der Mannigfaltigkeitslehre,''
\textit{Jahresbericht der Deutschen Mathematiker-Vereinigung} \textbf{1}, 75--78 (1891).

\bibitem{Bekenstein1981}
J.~D.~Bekenstein,
``Universal upper bound on the entropy-to-energy ratio for bounded systems,''
\textit{Phys.\ Rev.\ D} \textbf{23}, 287 (1981).

\bibitem{Horndeski1974}
G.~W.~Horndeski,
``Second-order scalar-tensor field equations in a four-dimensional space,''
\textit{Int.\ J.\ Theor.\ Phys.}\ \textbf{10}, 363 (1974).

\bibitem{Khoury2004}
J.~Khoury and A.~Weltman,
``Chameleon fields: Awaiting surprises for tests of gravity in space,''
\textit{Phys.\ Rev.\ Lett.}\ \textbf{93}, 171104 (2004).

\bibitem{Connes1996}
A.~Connes,
``Gravity coupled with matter and the foundation of non-commutative geometry,''
\textit{Commun.\ Math.\ Phys.}\ \textbf{182}, 155--176 (1996).

\bibitem{Kapitanov2025}
F.~Kapitanov,
``Ontological Resolution Theory: Gravity as Memory, Cognition as Indexing,''
Zenodo (2025). DOI: 10.5281/zenodo.17745212.

\end{thebibliography}

\end{document}